\RequirePackage[l2tabu, orthodox]{nag}
\pdfinfoomitdate=1
\pdfsuppressptexinfo=15
\pdftrailerid{}
\providecommand{\BotgeDarkMode}{1}
\documentclass[11pt,a4paper]{article}

\usepackage[T1]{fontenc}
\usepackage{lmodern}
\usepackage{microtype}
\usepackage[margin=1in]{geometry}
\usepackage{amsmath,amssymb,amsthm,mathtools}
\usepackage{booktabs}
\usepackage{xcolor}

\definecolor{BotgeBg}{HTML}{0e1117}
\definecolor{BotgeFg}{HTML}{e9edf4}
\definecolor{BotgeAccent}{HTML}{58a6ff}
\definecolor{BotgeLightAccent}{HTML}{0969da}

\usepackage{darkmode}
\ifnum\BotgeDarkMode=1\relax
  \enabledarkmode
  \pagecolor{BotgeBg}
  \color{BotgeFg}
  \colorlet{fgcolor}{BotgeFg}
\else
  \disabledarkmode
\fi

\usepackage{hyperref}
\hypersetup{%
  colorlinks=true,%
  bookmarksnumbered=true,%
  bookmarksopen=true,%
  pdfdisplaydoctitle=true,%
  pdfstartview=FitH,%
  pdftitle={A Bipartite Ising-Rx Circuit for Correlated Sampling},%
  pdfauthor={The Botge Project},%
  pdfsubject={Bipartite Ising-Rx Entanglement, Bitwise Correlations, and Rejection Sampling},%
  pdfkeywords={Ising coupling, IQP circuits, entanglement, mutual information, rejection sampling}%
}
\IfDarkModeTF{%
  \hypersetup{linkcolor=BotgeAccent, citecolor=BotgeAccent, urlcolor=BotgeAccent}%
}{%
  \hypersetup{linkcolor=BotgeLightAccent, citecolor=BotgeLightAccent, urlcolor=BotgeLightAccent}%
}
% Loaded before the theorem declarations so environments sharing a counter keep their own names.
\usepackage{aliascnt}
\usepackage[capitalise,noabbrev,nameinlink]{cleveref}

% One counter per section, shared by every numbered environment. Each environment counts with an
% alias of it, so cleveref names references by environment.
\newcommand{\SharedTheorem}[3]{%
  \newaliascnt{#1}{theorem}%
  \newtheorem{#1}[#1]{#2}%
  \aliascntresetthe{#1}%
  \crefname{#1}{#2}{#3}%
}
\theoremstyle{plain}
\newtheorem{theorem}{Theorem}[section]
\crefname{theorem}{Theorem}{Theorems}
\SharedTheorem{proposition}{Proposition}{Propositions}
\SharedTheorem{lemma}{Lemma}{Lemmas}
\SharedTheorem{corollary}{Corollary}{Corollaries}
\theoremstyle{definition}
\SharedTheorem{definition}{Definition}{Definitions}
\theoremstyle{remark}
\SharedTheorem{remark}{Remark}{Remarks}
\newtheorem*{remark*}{Remark}

\newcommand{\C}{\mathbb{C}}
\newcommand{\R}{\mathbb{R}}
\newcommand{\Z}{\mathbb{Z}}
\newcommand{\N}{\mathbb{N}}
\newcommand{\Hil}{\mathcal{H}}
\newcommand{\Un}{\mathrm{U}}
\newcommand{\Tr}{\operatorname{Tr}}
\newcommand{\Prob}{\mathbb{P}}
\newcommand{\Exp}{\mathbb{E}}
\newcommand{\Var}{\operatorname{Var}}
\newcommand{\Cov}{\operatorname{Cov}}
\newcommand{\Corr}{\operatorname{Corr}}
\newcommand{\ket}[1]{\lvert#1\rangle}
\newcommand{\bra}[1]{\langle#1\rvert}
\newcommand{\braket}[2]{\langle#1\vert#2\rangle}
\newcommand{\ketbra}[2]{\lvert#1\rangle\mkern-2mu\langle#2\rvert}
\newcommand{\proj}[1]{\ketbra{#1}{#1}}

\title{%
  A Bipartite Ising-\texorpdfstring{\( \mathrm{Rx} \)}{Rx} Circuit for Correlated Sampling:\\
  Entanglement, Bitwise Correlations, and Rejection Sampling%
}
\author{The Botge Project}
\date{}

\begin{document}
\maketitle

\begin{abstract}
  We study a quantum circuit on two registers, \( B \) and \( O \), of \( n_b \) and \( n_o \)
  qubits with \( n_b + n_o \le 16 \), and prove the properties that make its measurement a
  well-behaved correlated sampler. The circuit prepares a uniform superposition, couples the
  centered positions of the two registers through commuting \( \mathrm{Rzz} \) rotations, and
  reads out with the local gate \( \mathrm{Rx}(-\pi / 4) \), which we show is the optimal
  readout angle. The registers are always entangled but never maximally, every single measured
  bit is uniform, the most significant bits are positively correlated, and so the mutual
  information \( I(B; O) \) is positive. Rejection sampling onto any admissible domain accepts
  with probability above \( 1/4 \), and an interval-arithmetic computation sharpens these bounds
  across all \( 414 \) admissible configurations. The analytical results are formalized in
  Lean~4 with Mathlib.
\end{abstract}

\noindent\textbf{Keywords:} Ising coupling, IQP circuits, entanglement entropy, mutual
information, rejection sampling, interval arithmetic, formal verification.

\section{Introduction}\label{sec:introduction}

Given two finite index sets of sizes \( C_b \) and \( C_o \), we draw a pair \( (x_b, x_o) \)
whose components are correlated while each stays spread across its set, by measuring a
two-register state in the computational basis. We size the registers in~\cref{sec:budget},
define the circuit in~\cref{sec:circuit}, prove its entanglement and correlation properties
in~\cref{sec:entanglement}, and analyze rejection sampling onto the requested sets
in~\cref{sec:rejection}. The circuit is an instantaneous quantum polynomial-time (IQP)
circuit~\cite{ShepherdBremner2009,BremnerJozsaShepherd2011} whose single bilinear coupling makes
it exactly solvable, and the sampler is von Neumann's rejection method~\cite{vonNeumann1951}.

\paragraph{Design choices.} The registers share a budget of \( 16 \) qubits, split by the
caps~\eqref{eq:pool-cap} as \( 9 + 7 \) when \( C_b \le 512 \) and as \( 10 + 6 \) otherwise.
The coupling strength \( u_L = \varphi^{-1}\sqrt{L} \) grows with the depth \( L \) and stays
below \( \pi / 2 \) exactly for \( L \le 6 \), since
\( u_6 \approx 1.5139 < \pi / 2 \approx 1.5708 < u_7 \approx 1.6352 \). The algebraic factor
\( \varphi^{-1} \) rules out a vanishing correlation between the most significant bits
(\cref{thm:bitwise-corr}~(3)), the readout angle \( -\pi / 4 \) uniquely maximizes the
correlation (\cref{thm:bitwise-corr}~(2)), and the domain floor \( M^{\min}(n) \) yields
the acceptance floor of~\cref{thm:rejection}.

\paragraph{Verification.} All results are formalized in Lean~4~\cite{deMouraUllrich2021} with
Mathlib~\cite{Mathlib2020}, using no axioms beyond Mathlib's foundations; item~(3)
of~\cref{thm:bitwise-corr} is formalized on the admissible triples. The gate tally
of~\cref{prop:compile} and the passage from independent trials to the stopping-time series
in~\cref{thm:rejection} are argued on paper only. The bounds of~\cref{prop:certificate}
and~\cref{tab:domain-profile} are proved by exhaustive interval arithmetic~\cite{Moore1966},
and every other printed decimal is proved to round to its printed digits.

\subsection{Notation}\label{sec:notation}

\begin{itemize}
  \item \textbf{Registers.} For \( n \ge 1 \), \( N \coloneqq 2^n \) and
    \( \Hil_n \coloneqq {(\C^2)}^{\otimes n} \) with little-endian basis
    \( \ket{x} \coloneqq \ket{x_0} \otimes \cdots \otimes \ket{x_{n-1}} \),
    \( x = \sum_{j} x_j 2^j \in \Z_N \coloneqq \{0, \dots, N - 1\} \). Register \( B \) has
    \( n_b \) qubits and \( O \) has \( n_o \).
  \item \textbf{Gates.} \( X, Y, Z \) are the Pauli matrices,
    \( \ket{\pm} \coloneqq (\ket{0} \pm \ket{1}) / \sqrt{2} \),
    \( T \coloneqq \mathrm{diag}(1, e^{i\pi / 4}) \),
    \( \mathrm{Rx}(\alpha) \coloneqq \exp(-i\frac{\alpha}{2} X) \),
    \( \mathrm{Rz}(\alpha) \coloneqq \exp(-i\frac{\alpha}{2} Z) \) and
    \( \mathrm{Rzz}(\theta) \coloneqq \exp(-i\frac{\theta}{2} Z \otimes Z) \)~\cite{NielsenChuang}.
    \( A_q \) is \( A \) on qubit \( q \) and the identity elsewhere; \( A_{b,j} \) and
    \( A_{o,\ell} \) act on qubit \( j \) of \( B \) and \( \ell \) of \( O \), and
    \( \mathrm{Rzz}_{p,q} \) on qubits \( p, q \).
  \item \textbf{Measurement.} Measuring \( \ket{\psi} \in \Hil_{n_b} \otimes \Hil_{n_o} \) gives
    \( (B, O) \in \Z_{N_b} \times \Z_{N_o} \) with probability
    \( P(x_b, x_o) = {\lvert \braket{x_b, x_o}{\psi} \rvert}^2 \) and bits \( B_j, O_\ell \).
    For \( \ket{\psi_L} \) (\cref{def:ising-rx}) this law is \( P_L \), with marginals
    \( P_{B,L} \) and \( P_{O,L} \).
  \item \textbf{Entropies} are in bits with \( 0 \log 0 \coloneqq 0 \)~\cite{NielsenChuang}:
    \( S(\rho) \coloneqq -\Tr(\rho \log_2 \rho) \), \( S_2(\rho) \coloneqq -\log_2 \Tr(\rho^2) \),
    \( h_2(p) \coloneqq -p \log_2 p - (1 - p) \log_2(1 - p) \), the Shannon entropy
    \( H(\cdot) \), \( H(O \mid B) \coloneqq H(B, O) - H(B) \) and
    \( I(B; O) \coloneqq H(B) + H(O) - H(B, O) \); \( \Tr_O \) is the partial trace over \( O \).
  \item \textbf{Golden ratio.} \( \varphi^{-1} \coloneqq \frac{\sqrt{5} - 1}{2} \approx 0.618034 \).
\end{itemize}

\clearpage
\section{Register Sizing}\label{sec:budget}

\begin{proposition}[Register Widths and Occupancy]\label{prop:budget}
  Let \( C_b, C_o \ge 2 \) be integer domain sizes. Under the piecewise cap
  \begin{equation}\label{eq:pool-cap}
    (M_b^{\max}, M_o^{\max}) \coloneqq
    \begin{cases}
      (512, 128), & C_b \le 512, \\
      (1024, 64), & C_b > 512,
    \end{cases}
  \end{equation}
  let \( M_b = \min(C_b, M_b^{\max}) \) and \( M_o = \min(C_o, M_o^{\max}) \), and
  \begin{equation}\label{eq:bit-widths}
    n_b = \lceil \log_2 M_b \rceil, \qquad n_o = \lceil \log_2 M_o \rceil, \qquad
    N_b = 2^{n_b}, \qquad N_o = 2^{n_o}.
  \end{equation}
  Then:
  \begin{enumerate}
    \item \( n_b \in \{1, \dots, 10\} \), \( n_o \in \{1, \dots, 7\} \) and
      \( n_{\mathrm{tot}} \coloneqq n_b + n_o \le 16 \);
    \item \( N_b / 2 < M_b \le N_b \) and \( N_o / 2 < M_o \le N_o \), with \( M = N = 2 \) when
      \( n = 1 \);
    \item \( \dim(\Hil_{n_b} \otimes \Hil_{n_o}) = 2^{n_{\mathrm{tot}}} \le 2^{16} = 65{,}536 \).
  \end{enumerate}
\end{proposition}
\begin{proof}
  The two branches of~\eqref{eq:pool-cap} bound \( (M_b, M_o) \) by \( (2^9, 2^7) \) and
  \( (2^{10}, 2^6) \), giving~(1). Since \( M \ge 2 \), \( n \ge 1 \) and the ceiling gives
  \( 2^{n-1} < M \le 2^n \), which is~(2); (3) follows from~(1).
\end{proof}

\paragraph{Admissible configurations.} A triple \( (n_b, n_o, L) \) is \emph{admissible} when
\( n_b \in \{1, \dots, 10\} \), \( n_o \in \{1, \dots, 7\} \), \( n_b + n_o \le 16 \) and
\( L \in \{1, \dots, 6\} \); these are the widths of~\cref{prop:budget} at every depth,
\( 69 \times 6 = 414 \) triples. With \( M^{\min}(n) \coloneqq \max(2, 2^{n-1} + 1) \), domain
sizes are \emph{admissible} when \( M^{\min}(n_b) \le M_b \le N_b \) and
\( M^{\min}(n_o) \le M_o \le N_o \).

\clearpage
\section{The Circuit}\label{sec:circuit}

\begin{definition}[Centered Position Coupling and Ising-\( \mathrm{Rx} \) State]\label{def:ising-rx}
  On an \( n \)-qubit register, let
  \( \mathcal{M}_N \coloneqq \{-\tfrac{N-1}{2}, -\tfrac{N-3}{2}, \dots, +\tfrac{N-1}{2}\} \) and
  \begin{equation}\label{eq:centered-x}
    \hat{\tilde{x}} \coloneqq \sum_{x=0}^{N-1} \left(x - \frac{N - 1}{2}\right)\proj{x},
  \end{equation}
  whose eigenvalues range over \( \mathcal{M}_N \). For depth \( L \in \{1, \dots, 6\} \), let
  \( u_L \coloneqq \varphi^{-1}\sqrt{L} \), so that \( 0.6180 \le u_L \le 1.5139 < \pi / 2 \), and
  \begin{align}
    \hat{V}_L &\coloneqq -\frac{2 u_L}{\max(N_b, N_o)}\, \hat{\tilde{x}}_b \otimes \hat{\tilde{x}}_o,
    \label{eq:pos-ham} \\
    R_X &\coloneqq \mathrm{Rx}\left(-\frac{\pi}{4}\right) = \exp\left(i\frac{\pi}{8} X\right) =
    H\,\mathrm{Rz}\left(-\frac{\pi}{4}\right) H = e^{i\pi / 8} H T^\dagger H, \label{eq:rx-gate} \\
    U_L &\coloneqq \left(R_X^{\otimes n_b} \otimes R_X^{\otimes n_o}\right) e^{-i\hat{V}_L}
    \left(H^{\otimes n_b} \otimes H^{\otimes n_o}\right), \label{eq:ising-rx-unitary} \\
    \ket{\psi_L} &\coloneqq U_L \ket{0}^{\otimes(n_b + n_o)}.
    \label{eq:ising-rx-state}
  \end{align}
\end{definition}

\begin{proposition}[Gate Decomposition]\label{prop:compile}
  With no global phase, \( U_L \) is the product of \( n_b n_o + 2(n_b + n_o) \le 96 \) gates:
  \( n_b + n_o \) Hadamards, then the \( n_b n_o \) pairwise-commuting gates
  \[
    e^{-i\hat{V}_L} = \prod_{j=0}^{n_b - 1}\prod_{\ell=0}^{n_o - 1}
    \mathrm{Rzz}_{b_j, o_\ell}(\theta_{j\ell}),
  \]
  with
  \begin{equation}\label{eq:rzz-pos}
    \theta_{j\ell} = -\frac{u_L\, 2^j 2^\ell}{\max(N_b, N_o)} < 0,
  \end{equation}
  then \( n_b + n_o \) gates \( \mathrm{Rx}(-\pi / 4) \).
\end{proposition}
\begin{proof}
  Substituting \( x_j = \frac{1}{2}(I - Z_j) \) into \( x = \sum_j 2^j x_j \) and using
  \( \sum_j 2^j = N - 1 \),
  \begin{equation}\label{eq:pauli-x}
    \hat{\tilde{x}} = -\frac{N - 1}{2} I + \sum_{j=0}^{n-1} 2^j \frac{I - Z_j}{2} =
    -\frac{1}{2}\sum_{j=0}^{n-1} 2^j Z_j,
  \end{equation}
  so \( \hat{V}_L = \frac{1}{2}\sum_{j,\ell} \theta_{j\ell} Z_{b,j} Z_{o,\ell} \). These terms
  commute, which gives the product exactly; \( n_b n_o \le {(n_b + n_o)}^2 / 4 \le 64 \).
\end{proof}

\clearpage
\section{Entanglement and Correlations}\label{sec:entanglement}

Throughout, \( N_{\min} \coloneqq \min(N_b, N_o) \), \( N_{\max} \coloneqq \max(N_b, N_o) \),
and \( \ket{\psi_L^{(0)}} \coloneqq e^{-i\hat{V}_L}\ket{+}^{\otimes (n_b + n_o)} \) is the state
before readout, so \( \ket{\psi_L} = (R_X^{\otimes n_b} \otimes R_X^{\otimes n_o})\ket{\psi_L^{(0)}} \).

\begin{theorem}[Toeplitz Reduced State and Entanglement Bounds]\label{thm:entanglement}
  Let \( \rho_B \coloneqq \Tr_O(\proj{\psi_L}) \), \( \rho_O \coloneqq \Tr_B(\proj{\psi_L}) \),
  let \( R \in \{B, O\} \) have \( N_R = N_{\min} \) and \( \bar{R} \) be the other register, and
  \( \rho_R^{(0)} \coloneqq \Tr_{\bar{R}}(\proj{\psi_L^{(0)}}) \). Then:
  \begin{enumerate}
    \item \( \rho_R = W \rho_R^{(0)} W^\dagger \) with \( W \coloneqq R_X^{\otimes n_R} \), where
      \( \rho_R^{(0)} \) is the real symmetric Toeplitz matrix
      \( \bra{x_R}\rho_R^{(0)}\ket{x_R'} = \frac{1}{N_{\min}} g_L(x_R - x_R') \) with
      \( g_L(0) = 1 \) and \( g_L(\Delta) \coloneqq \frac{\sin(u_L \Delta)}{N_{\max}
      \sin(u_L \Delta / N_{\max})} \) otherwise, and
      \begin{equation}\label{eq:purity}
        \Tr(\rho_B^2) = \Tr(\rho_O^2) = \frac{1}{N_{\min}} +
        \frac{2}{N_{\min}^2}\sum_{\Delta=1}^{N_{\min} - 1} (N_{\min} - \Delta)\, {g_L(\Delta)}^2.
      \end{equation}
    \item The purity satisfies
      \begin{equation}\label{eq:purity-bounds}
        \frac{1}{N_{\min}} + \frac{2(N_{\min} - 1)}{N_{\min}^2}{\left(\frac{\sin u_L}{u_L}\right)}^2
        < \Tr(\rho_B^2) \le
        1 - \frac{2(N_{\min} - 1)}{N_{\min}^2}\sin^2\left(\frac{u_L}{2}\right) < 1,
      \end{equation}
      so \( 0 < S_2(\rho_B) \le S(\rho_B) < \min(n_b, n_o) \).
  \end{enumerate}
\end{theorem}
\begin{proof}
  The readout is local, so \( \rho_R = W \rho_R^{(0)} W^\dagger \). With
  \( \gamma_L = 2 u_L / N_{\max} \),
  \[
    \ket{\psi_L^{(0)}} = \frac{1}{\sqrt{N_b N_o}} \sum_{x_b, x_o}
    e^{i\gamma_L \tilde{x}_b \tilde{x}_o} \ket{x_b, x_o},
  \]
  and summing the geometric series over the consecutive integers
  \( \tilde{x}_{\bar{R}} \in \mathcal{M}_{N_{\max}} \) gives, for \( \Delta = x_R - x_R' \),
  \[
    \bra{x_R}\rho_R^{(0)}\ket{x_R'} = \frac{1}{N_{\min} N_{\max}}
    \sum_{\tilde{x}_{\bar{R}} \in \mathcal{M}_{N_{\max}}} e^{i \gamma_L \Delta \tilde{x}_{\bar{R}}} =
    \frac{1}{N_{\min}} g_L(\Delta),
  \]
  whose denominator is nonzero because \( 0 < u_L \lvert \Delta \rvert / N_{\max} < \pi / 2 \).
  Grouping \( {\lvert \bra{x_R}\rho_R^{(0)}\ket{x_R'} \rvert}^2 \) by \( \Delta \) and the Schmidt
  equality \( \Tr(\rho_B^2) = \Tr(\rho_O^2) \)~\cite{NielsenChuang} give~\eqref{eq:purity}.

  Since \( t \mapsto \sin t / t \) decreases on \( (0, \pi / 2) \) and
  \( u_L / N_{\max} \le u_L / 2 < \pi / 4 \),
  \( 2\sin(u_L / 2) \le N_{\max} \sin(u_L / N_{\max}) < u_L \); dividing
  \( \sin u_L = 2\sin(u_L / 2)\cos(u_L / 2) \) by these yields
  \[
    0.659 < \frac{\sin u_L}{u_L} < g_L(1) \le \cos\left(\frac{u_L}{2}\right) \le
    \cos\left(\frac{\varphi^{-1}}{2}\right) < 0.953.
  \]
  Dropping the \( \Delta \ge 2 \) terms of~\eqref{eq:purity} gives the lower bound;
  \( {g_L(\Delta)}^2 \le 1 \) (from \( \lvert \sin(N_{\max} z) \rvert \le N_{\max}
  \lvert \sin z \rvert \)) and \( {g_L(1)}^2 \le 1 - \sin^2(u_L / 2) \) give the upper bound,
  because the weights in~\eqref{eq:purity} sum to \( 1 \). Finally \( \rho_B \) has at most
  \( N_{\min} \) nonzero eigenvalues and purity above \( 1 / N_{\min} \), so
  \( S(\rho_B) < \log_2 N_{\min} = \min(n_b, n_o) \).
\end{proof}

\begin{lemma}[Trigonometric and Heisenberg Identities]\label{lem:identities}
  Let \( \hat{\Theta}_{b,j} \coloneqq \sum_{\ell} \theta_{j\ell} Z_{o,\ell} \) and
  \( \hat{\Theta}_{o,\ell} \coloneqq \sum_{j} \theta_{j\ell} Z_{b,j} \). Then:
  \begin{enumerate}
    \item (Morrie) for \( m \ge 0 \) and \( \sin x \neq 0 \),
      \begin{equation}\label{eq:morrie}
        \prod_{k=0}^{m-1} \cos(2^k x) = \frac{\sin(2^m x)}{2^m \sin x};
      \end{equation}
    \item for every angle \( \alpha \),
      \begin{equation}\label{eq:heis-rx}
        {\mathrm{Rx}(\alpha)}^\dagger Z\, \mathrm{Rx}(\alpha) = \cos(\alpha)\, Z + \sin(\alpha)\, Y,
      \end{equation}
      so \( R_X^\dagger Z R_X = \frac{1}{\sqrt{2}}(Z - Y) \);
    \item
      \begin{equation}\label{eq:heis-vl}
        e^{+i\hat{V}_L} Y_{b,j} e^{-i\hat{V}_L} = Y_{b,j}\cos\hat{\Theta}_{b,j} +
        X_{b,j}\sin\hat{\Theta}_{b,j},
      \end{equation}
      and likewise for \( Y_{o,\ell} \);
    \item for distinct qubits \( S \) and \( \beta_q \in \R \),
      \begin{equation}\label{eq:trig-plus-exp}
        \bra{+}^{\otimes \lvert S \rvert} \cos\left(\sum_{q \in S} \beta_q Z_q\right)
        \ket{+}^{\otimes \lvert S \rvert} = \prod_{q \in S} \cos\beta_q, \qquad
        \bra{+}^{\otimes \lvert S \rvert} \sin\left(\sum_{q \in S} \beta_q Z_q\right)
        \ket{+}^{\otimes \lvert S \rvert} = 0.
      \end{equation}
  \end{enumerate}
\end{lemma}
\begin{proof}
  (1)~Apply \( 2\sin y\cos y = \sin(2y) \) \( m \) times.
  (2)~\( X \) anticommutes with \( Z \), so the product equals
  \( Z e^{-i\alpha X} = \cos(\alpha)\, Z - i\sin(\alpha)\, ZX \), and \( ZX = iY \).
  (3)~By~\eqref{eq:pauli-x}, the terms of \( \hat{V}_L \) on \( b_j \) sum to
  \( A = \frac{1}{2} Z_{b,j} \hat{\Theta}_{b,j} \), which anticommutes with \( Y_{b,j} \); the
  rest commute, so the conjugate is \( Y_{b,j} e^{-2iA} \), and \( YZ = iX \).
  (4)~\( \bra{+}^{\otimes \lvert S \rvert} e^{i\sum_q \beta_q Z_q} \ket{+}^{\otimes \lvert S \rvert}
  = \prod_q \cos\beta_q \) is real; take real and imaginary parts.
\end{proof}

\begin{theorem}[Symmetry and Bitwise Correlations]\label{thm:bitwise-corr}
  Let \( Z_{b,j}, Z_{o,\ell} \in \{\pm 1\} \) be the observables of bits \( B_j \) and
  \( O_\ell \) (\( +1 \) on bit~\( 0 \)). Then:
  \begin{enumerate}
    \item \textbf{Symmetry and uniform bits.} \( X^{\otimes(n_b + n_o)}\ket{\psi_L} = \ket{\psi_L} \),
      so \( P_L(N_b - 1 - x_b,\, N_o - 1 - x_o) = P_L(x_b, x_o) \), both marginals are
      palindromic, and \( \Exp[Q] = 0 \) for every odd-weight product \( Q \) of \( Z_{b,j} \)
      and \( Z_{o,\ell} \). Hence every bit is uniform, \( \Prob(B < N_b / 2) = \Prob(O < N_o / 2)
      = \frac{1}{2} \), and \( \Prob(B < M_b), \Prob(O < M_o) \ge \frac{1}{2} \) for all admissible
      \( (M_b, M_o) \).
    \item \textbf{Bitwise correlations and optimal readout.} With readout
      \( {\mathrm{Rx}(\alpha)}^{\otimes(n_b + n_o)} \) in place of \( R_X^{\otimes(n_b + n_o)} \),
      \begin{equation}\label{eq:bitwise-corr}
        \Corr_\alpha(Z_{b,j}, Z_{o,\ell}) = -\frac{\sin(2\alpha)}{2}\sin\lvert\theta_{j\ell}\rvert
        \left(\prod_{\substack{j'=0 \\ j' \neq j}}^{n_b - 1} \cos\theta_{j'\ell} +
        \prod_{\substack{\ell'=0 \\ \ell' \neq \ell}}^{n_o - 1} \cos\theta_{j\ell'}\right),
      \end{equation}
      whose prefactor \( -\sin(2\alpha) \) is maximal on \( [-\pi / 2, \pi / 2] \) only at
      \( \alpha = -\pi / 4 \). At \( \alpha = -\pi / 4 \), for \( 0 \le j < j' < n_b \),
      \begin{equation}\label{eq:intra-corr}
        \Corr(Z_{b,j}, Z_{b,j'}) = \frac{1}{4}\left(
          \frac{\sin\left(\frac{u_L N_o (2^{j'} - 2^j)}{N_{\max}}\right)}%
               {N_o \sin\left(\frac{u_L (2^{j'} - 2^j)}{N_{\max}}\right)} -
          \frac{\sin\left(\frac{u_L N_o (2^{j'} + 2^j)}{N_{\max}}\right)}%
               {N_o \sin\left(\frac{u_L (2^{j'} + 2^j)}{N_{\max}}\right)}
        \right),
      \end{equation}
      with denominator angles in \( (0, 3\pi / 8) \), and symmetrically for \( O \).
    \item \textbf{Most significant bits and mutual information.} At \( \alpha = -\pi / 4 \),
      \begin{equation}\label{eq:msb-corr}
        \Corr(Z_{b, n_b - 1}, Z_{o, n_o - 1}) =
        \sin^2\left(\frac{u_L N_{\min}}{4}\right)
        \left(\frac{1}{N_{\max}\sin\left(\frac{u_L N_{\min}}{2 N_{\max}}\right)} +
        \frac{1}{N_{\min}\sin\left(\frac{u_L}{2}\right)}\right),
      \end{equation}
      which is positive for all \( n_b, n_o \ge 1 \) and \( L \le 6 \) by the
      Lindemann--Weierstrass theorem, so
      \begin{equation}\label{eq:mi-lower}
        I(B; O) \ge I(B_{n_b - 1}; O_{n_o - 1}) =
        1 - h_2\left(\frac{1 + \Corr(Z_{b, n_b - 1}, Z_{o, n_o - 1})}{2}\right) > 0.
      \end{equation}
  \end{enumerate}
\end{theorem}
\begin{proof}
  \textbf{(1)} Let \( \mathcal{X} \coloneqq X^{\otimes(n_b + n_o)} \). By~\eqref{eq:pauli-x},
  \( \mathcal{X} \) flips the sign of both \( \hat{\tilde{x}}_b \) and \( \hat{\tilde{x}}_o \), so
  it commutes with \( e^{-i\hat{V}_L} \), with every readout layer, and fixes
  \( \ket{+}^{\otimes(n_b + n_o)} \); hence \( \mathcal{X}\ket{\psi_L} = \ket{\psi_L} \). Since
  \( X^{\otimes n}\ket{x} = \ket{N - 1 - x} \), this is the complement symmetry, and
  \( \mathcal{X} Q \mathcal{X} = -Q \) forces \( \Exp[Q] = 0 \). As
  \( \{B < N_b / 2\} = \{B_{n_b - 1} = 0\} \) and \( M_b \ge N_b / 2 \), the bounds follow.

  \textbf{(2)} The argument of~(1) holds for every readout angle, so each \( Z \) has mean
  \( 0 \) and variance \( 1 \) and \( \Corr_\alpha = \Exp_\alpha[Z_{b,j} Z_{o,\ell}] \).
  By~\eqref{eq:heis-rx}, this is the expectation in \( \ket{+}^{\otimes(n_b + n_o)} \),
  conjugated by \( e^{-i\hat{V}_L} \), of
  \( \cos^2(\alpha)\, Z_{b,j} Z_{o,\ell} + \sin^2(\alpha)\, Y_{b,j} Y_{o,\ell} +
  \frac{\sin 2\alpha}{2}(Z_{b,j} Y_{o,\ell} + Y_{b,j} Z_{o,\ell}) \). The \( ZZ \) term
  commutes with \( \hat{V}_L \) and vanishes. In \( Z_{b,j} Y_{o,\ell} \),
  identity~\eqref{eq:heis-vl} leaves \( Z_{b,j} X_{o,\ell} \sin\hat{\Theta}_{o,\ell} \); splitting off
  \( \theta_{j\ell} Z_{b,j} \) from \( \hat{\Theta}_{o,\ell} \) and
  applying~\eqref{eq:trig-plus-exp} gives \( -\sin\lvert\theta_{j\ell}\rvert
  \prod_{j' \neq j} \cos\theta_{j'\ell} \), and symmetrically for \( Y_{b,j} Z_{o,\ell} \). In
  \( Y_{b,j} Y_{o,\ell} \), the shared term \( \theta_{j\ell} Z_{b,j} Z_{o,\ell} \) and all terms
  off~\( \{ b_j, o_\ell \} \) commute with it, and every surviving term carries a \( Y \) or a
  \( \sin \) that vanishes by~\eqref{eq:trig-plus-exp}. This gives~\eqref{eq:bitwise-corr}.
  For~\eqref{eq:intra-corr}, \( \hat{\Theta}_{b,j} \) and \( \hat{\Theta}_{b,j'} \) act only on
  \( O \), so only \( \frac{1}{2} X_{b,j} X_{b,j'} \sin\hat{\Theta}_{b,j}\sin\hat{\Theta}_{b,j'} \)
  survives; \( 2\sin A\sin B = \cos(A - B) - \cos(A + B) \), then~\eqref{eq:trig-plus-exp}
  and~\eqref{eq:morrie} with \( m = n_o \) finish, and
  \( 1 \le 2^{j'} \pm 2^j \le \frac{3}{4} N_{\max} \) keeps the angles in \( (0, 3\pi / 8) \).

  \textbf{(3)} Let \( x_b \coloneqq \frac{u_L N_b}{2 N_{\max}} \) and
  \( x_o \coloneqq \frac{u_L N_o}{2 N_{\max}} \), both in \( (0, \pi / 4) \). At
  \( (n_b - 1, n_o - 1) \), \( \lvert\theta\rvert = u_L N_{\min} / 4 \), and the remaining angles
  are \( 2^{j'} x_o \) and \( 2^{\ell'} x_b \), so~\eqref{eq:morrie} turns~\eqref{eq:bitwise-corr}
  into \( \sin^2(u_L N_{\min} / 4)\,(1 / (N_b \sin x_o) + 1 / (N_o \sin x_b)) \),
  which is~\eqref{eq:msb-corr}. Its first factor vanishes only if
  \( \pi = \varphi^{-1}\sqrt{L}\, N_{\min} / (4k) \) for an integer \( k \ge 1 \), impossible since
  \( \pi \) is transcendental. Writing \( c \) for this correlation, \( 0 < c \le 1 \), the two most significant
  bits have law \( \frac{1 \pm c}{4} \), so \( I(B_{n_b - 1}; O_{n_o - 1}) = 1 - h_2(\frac{1 + c}{2})
  > 0 \), and the data processing inequality~\cite{NielsenChuang} gives~\eqref{eq:mi-lower}.
\end{proof}

\begin{remark*}[Item~(3) Without Transcendence]
  On the admissible triples, \( N_{\min} = 2^m \) with \( m \le 7 \), and for each of the
  \( 42 \) pairs \( (m, L) \), rational enclosures of \( \sqrt{5} \), \( \sqrt{L} \) and \( \pi \)
  place \( u_L 2^m / 4 \) strictly between consecutive multiples of \( \pi \). The Lean
  formalization proves item~(3) on these triples this way, without additional axioms.
\end{remark*}

\begin{corollary}[One Qubit per Register]\label{cor:two-by-two}
  For \( n_b = n_o = 1 \), let \( \vartheta_L \coloneqq u_L / 2 \in (0, \pi / 4) \). Then
  \begin{equation}\label{eq:two-by-two-state}
    \ket{\psi_L} = e^{i\pi / 4}
    \left(\cos\left(\frac{\vartheta_L}{2}\right)\ket{++} +
    \sin\left(\frac{\vartheta_L}{2}\right)\ket{--}\right),
  \end{equation}
  with \( P_L(0, 0) = P_L(1, 1) = \frac{1 + \sin\vartheta_L}{4} \) and
  \( P_L(0, 1) = P_L(1, 0) = \frac{1 - \sin\vartheta_L}{4} \). In particular:
  \begin{enumerate}
    \item \( S(\rho_B) = h_2\left(\frac{1 + \cos\vartheta_L}{2}\right) > 0 \)
      (\( \approx 0.3656 \) bits at \( L = 3 \));
    \item each register is uniform;
    \item \( \Corr(Z_b, Z_o) = \sin\vartheta_L \in [0.304, 0.687] \), strictly increasing in
      \( L \), and \( I(B; O) = 1 - h_2\left(\frac{1 + \sin\vartheta_L}{2}\right) > 0 \)
      (\( 0.1968 \) bits at \( L = 3 \), \( 0.3738 \) bits at \( L = 6 \)).
  \end{enumerate}
\end{corollary}
\begin{proof}
  By~\eqref{eq:rzz-pos} the coupling is \( \mathrm{Rzz}(-\vartheta_L) \), mapping
  \( \ket{++} \mapsto \cos(\vartheta_L / 2)\ket{++} + i\sin(\vartheta_L / 2)\ket{--} \). As
  \( R_X\ket{\pm} = e^{\pm i\pi / 8}\ket{\pm} \) and \( i e^{-i\pi / 4} = e^{i\pi / 4} \), the
  readout gives~\eqref{eq:two-by-two-state}, whose amplitudes
  \( \frac{e^{i\pi / 4}}{2}(\cos\frac{\vartheta_L}{2} \pm \sin\frac{\vartheta_L}{2}) \) have the
  stated probabilities. Then \( \rho_B = \frac{1}{2}(I + \cos\vartheta_L X) \) gives~(1), and
  then~\cref{thm:bitwise-corr} with \( N_{\min} = N_{\max} = 2 \) gives~(2),
  \( \Corr = \sin^2\vartheta_L / \sin\vartheta_L \), and equality in~\eqref{eq:mi-lower}.
\end{proof}

\begin{remark}[Classical Center Sampling Destroys Correlation]\label{rem:washout}
  Applying a translation-covariant unitary \( W \) to independent uniform basis states over
  \( \Z_{N_b} \times \Z_{N_o} \) and discarding them gives
  \[
    \Prob(B = x_b, O = x_o) = \frac{1}{N_b N_o}\sum_{x_{b0}, x_{o0}}
    {\lvert \bra{x_b - x_{b0}, x_o - x_{o0}} W \ket{0, 0} \rvert}^2 = \frac{1}{N_b N_o},
  \]
  so \( I(B; O) = 0 \); preparing \( \ket{\psi_L} \) directly keeps the
  positive mutual information of~\cref{thm:bitwise-corr}.
\end{remark}

\clearpage
\section{Rejection Sampling}\label{sec:rejection}

Fix an admissible triple and admissible \( (M_b, M_o) \), let
\( A \coloneqq \{0, \dots, M_b - 1\} \times \{0, \dots, M_o - 1\} \), and let
\( X_1, X_2, \dots \) be independent measurements of \( \ket{\psi_L} \), each with law \( P_L \).

\begin{theorem}[Rejection Sampling]\label{thm:rejection}
  Let \( p_A(M_b, M_o, L) \coloneqq \sum_{x \in A} P_L(x) \) and
  \( \tau \coloneqq \inf \{r \ge 1 : X_r \in A\} \). Then:
  \begin{enumerate}
    \item \textbf{Monotonicity.}
      \begin{equation}\label{eq:monotonicity}
        p_A(M_b, M_o, L) \ge p_A\big(M^{\min}(n_b), M^{\min}(n_o), L\big),
      \end{equation}
      and \( p_A(N_b, N_o, L) = 1 \).
    \item \textbf{Acceptance floor.} With \( c_{n_b, n_o, L} \coloneqq
      \Corr(Z_{b, n_b - 1}, Z_{o, n_o - 1}) > 0 \) from~\eqref{eq:msb-corr},
      \begin{equation}\label{eq:analytic-floor}
        p_A(M_b, M_o, L) \ge \frac{1 + c_{n_b, n_o, L}}{4} > \frac{1}{4},
      \end{equation}
      with \( p_A \ge \frac{1}{2} \) when \( \min(n_b, n_o) = 1 \) and \( p_A = 1 \) when
      \( n_b = n_o = 1 \).
    \item \textbf{Termination and accepted law.} \( \Prob(\tau < \infty) = 1 \),
      \( \Exp[\tau] = 1 / p_A < 4 \),
      \( \Prob(\tau > 256) = {(1 - p_A)}^{256} < {\left(\frac{3}{4}\right)}^{256} \approx
      1.04 \times 10^{-32} \), and
      \begin{equation}\label{eq:cond-law}
        \Prob(X_\tau = x) = \frac{P_L(x)}{p_A(M_b, M_o, L)} \qquad \forall x \in A.
      \end{equation}
  \end{enumerate}
\end{theorem}
\begin{proof}
  (1)~\( A \) contains the rectangle of the smallest admissible domain, and probabilities are
  non-negative; unitarity gives \( p_A(N_b, N_o, L) = 1 \).
  (2)~If \( n_b = n_o = 1 \), \( A \) is everything. If exactly one width is \( 1 \), that register
  is fully accepted and the other accepts more than its lower half, which has probability
  \( \frac{1}{2} \) by~\cref{thm:bitwise-corr}~(1). Otherwise \( A \) contains the quadrant
  \( \{B_{n_b - 1} = 0, O_{n_o - 1} = 0\} \), whose probability is
  \[
    \frac{1 + \Exp[Z_{b, n_b - 1}] + \Exp[Z_{o, n_o - 1}] +
    \Exp[Z_{b, n_b - 1} Z_{o, n_o - 1}]}{4} = \frac{1 + c_{n_b, n_o, L}}{4} > \frac{1}{4}.
  \]
  (3)~The trials are independent with success probability \( p_A > \frac{1}{4} \), so
  \( \Prob(\tau > k) = {(1 - p_A)}^k \), \( \Exp[\tau] = 1 / p_A \), and
  \[
    \Prob(X_\tau = x) = \sum_{k=1}^\infty {(1 - p_A)}^{k-1} P_L(x) =
    \frac{P_L(x)}{p_A(M_b, M_o, L)}. \qedhere
  \]
\end{proof}

\begin{table}[htbp]
  \centering
  \small
  \caption{Metrics of representative admissible triples, certified to the last printed
    digit.\label{tab:domain-profile}}
  \begin{tabular}{@{}lrrrrrr@{}}
    \toprule
    \((n_b, n_o, L)\) & \(S_2(\rho_B)\) [bits] & \(p_A^{\min}\) & \(\min H/n\) & Max/min ratio & \(I(B; O)\) [bits] & \(\min P_L\) \\
    \midrule
    \((1, 1, 3)\)   & \(0.2010\) & \(1.0000\) & \(1.0000\) & \(1.00\) & \(0.1968\) & \(1.22 \times 10^{-1}\) \\
    \((3, 2, 3)\)   & \(0.8807\) & \(0.5561\) & \(0.9838\) & \(1.87\) & \(0.6752\) & \(8.65 \times 10^{-4}\) \\
    \((5, 4, 3)\)   & \(2.5835\) & \(0.3992\) & \(0.9888\) & \(2.22\) & \(0.7369\) & \(2.24 \times 10^{-6}\) \\
    \((8, 6, 3)\)   & \(4.4895\) & \(0.2861\) & \(0.9924\) & \(2.73\) & \(0.8828\) & \(3.51 \times 10^{-9}\) \\
    \((9, 7, 3)\)   & \(5.4703\) & \(0.2614\) & \(0.9933\) & \(2.88\) & \(0.9513\) & \(4.21 \times 10^{-10}\) \\
    \((10, 6, 3)\)  & \(4.4896\) & \(0.2821\) & \(0.9926\) & \(3.02\) & \(0.8833\) & \(1.37 \times 10^{-11}\) \\
    \((9, 7, 5)\)   & \(5.8337\) & \(0.2595\) & \(0.9921\) & \(3.23\) & \(1.1085\) & \(2.50 \times 10^{-10}\) \\
    \((9, 7, 6)\)   & \(5.9637\) & \(0.2638\) & \(0.9919\) & \(3.22\) & \(1.8157\) & \(6.25 \times 10^{-11}\) \\
    \bottomrule
  \end{tabular}
\end{table}

\begin{proposition}[Certified Bounds on All Admissible Triples]\label{prop:certificate}
  Evaluating~\eqref{eq:ising-rx-unitary} in rigorous interval arithmetic on all \( 414 \)
  admissible triples \( (n_b, n_o, L) \) gives the following bounds, where \( p_A \) is the
  acceptance probability on the smallest admissible domain \( (M^{\min}(n_b), M^{\min}(n_o)) \)
  (at \( (9, 7, 5) \), \( M_b = 257 \) and \( M_o = 65 \)):
  \begin{center}
    \small\upshape
    \begin{tabular}{@{}p{0.37\linewidth}p{0.37\linewidth}p{0.2\linewidth}@{}}
      \toprule
      Quantity & Value or bound & Attained at \\
      \midrule
      \( \min p_A \) & \( \approx 0.259488 \) & \( (9, 7, 5) \) \\
      \( \min \Prob(B < N_b / 2, O < N_o / 2) \) & \( \approx 0.250170 \) & \( (10, 6, 5) \) \\
      \( \max \Exp[\tau] = \max 1 / p_A \) & \( \le 3.854 \) & \( (9, 7, 5) \) \\
      \( \max \Prob(\tau > 256) \) & \( \le {(1 - \min p_A)}^{256} \approx 3.98 \times 10^{-34} \) & \( (9, 7, 5) \) \\
      \midrule
      \( \min \sum_{b < M^{\min}(n_b)} P_{B,L}(b) \) & \( \ge 0.500735 \) & \( (10, 2, 3) \) \\
      \( \min \sum_{o < M^{\min}(n_o)} P_{O,L}(o) \) & \( \ge 0.500735 \) & \\
      \( \min \min(H(B) / n_b, H(O) / n_o) \) & \( \ge 0.9785 \) & \( (2, 2, 5) \) \\
      \quad on \( 16 \)-qubit triples & \( \ge 0.9897 \) & \\
      \( \max \) marginal max/min ratio & \( \le 3.693 \) & \( (10, 4, 6) \) \\
      \midrule
      \( \min I(B; O) \) & \( \ge 0.06778 \) bits & \( (1, 1, 1) \) \\
      \( I(B; O) \) & \( \approx 0.9513 \) and \( 1.8157 \) bits & \( (9, 7, 3) \) and \( (9, 7, 6) \) \\
      \( \min \min(H(O \mid B) / n_o, H(B \mid O) / n_b) \) & \( \ge 0.5701 \) & \( (10, 1, 6) \) \\
      median of \( 2^{H(O \mid B = b)} \) over \( b \) & \( \approx 66.7 \) of \( 128 \) & \( (9, 7, 3) \) \\
      \midrule
      \( \min_{x_b, x_o} P_L(x_b, x_o) \) & \( \ge 1.366 \times 10^{-11} \) & \( (10, 6, 3) \) \\
      \bottomrule
    \end{tabular}
  \end{center}
\end{proposition}
\begin{proof}
  Every quantity is computed as an interval containing its exact value, with outward rounding
  and series for \( \sin \), \( \cos \), \( \log_2 \) and \( 2^x \) whose remainders are
  bounded. An inequality holds when the whole enclosure satisfies it, each \( \approx \) value is
  enclosed to within one unit of its last digit, and each named extremizer is unique because its
  enclosure lies strictly beyond all others.
\end{proof}

\clearpage
\section{Conclusion}\label{sec:conclusion}

A single bilinear coupling between uniformly prepared registers, read out with
\( \mathrm{Rx}(-\pi / 4) \), yields entangled registers with uniform bits and positive mutual
information, and rejection sampling onto them accepts with probability above \( 1/4 \). Extending
the formalization of~\cref{thm:bitwise-corr}~(3) beyond the admissible triples would need a
formal transcendence result for \( \pi \).

\clearpage
\begin{thebibliography}{9}

  \bibitem{BremnerJozsaShepherd2011}
  M.~J.~Bremner, R.~Jozsa, and D.~J.~Shepherd,
  Classical simulation of commuting quantum computations implies collapse of the polynomial
  hierarchy, \emph{Proc.\ R.\ Soc.\ A} \textbf{467} (2011), 459--472.

  \bibitem{deMouraUllrich2021}
  L.~de~Moura and S.~Ullrich,
  The Lean~4 theorem prover and programming language, \emph{CADE~28}, LNCS \textbf{12699}
  (2021), 625--635.

  \bibitem{Mathlib2020}
  The mathlib Community,
  The Lean mathematical library, \emph{CPP~2020}, ACM (2020), 367--381.

  \bibitem{Moore1966}
  R.~E.~Moore, \emph{Interval Analysis}, Prentice-Hall, 1966.

  \bibitem{NielsenChuang}
  M.~A.~Nielsen and I.~L.~Chuang,
  \emph{Quantum Computation and Quantum Information}, Cambridge University Press, 2010.

  \bibitem{ShepherdBremner2009}
  D.~Shepherd and M.~J.~Bremner,
  Temporally unstructured quantum computation, \emph{Proc.\ R.\ Soc.\ A} \textbf{465} (2009),
  1413--1439.

  \bibitem{vonNeumann1951}
  J.~von Neumann,
  Various techniques used in connection with random digits, \emph{NBS Appl.\ Math.\ Ser.}
  \textbf{12} (1951), 36--38.

\end{thebibliography}

\end{document}
